\clearpage\SourceLeaf{358}

\begin{center}
\large\bfseries [DIE KUGEL.]
\end{center}

\begin{center}
[1.]
\end{center}

\begin{center}
\bfseries ORTHOGRAFISCHE PROJEKTION DER KUGEL.
\end{center}

\begin{center}
Verändertes Aussehen eines unendlich kleinen Teils der Kugeloberfläche,\\
wenn die Projektionsebene geändert wird.
\end{center}

Man wählt als $x$-Achse die Schnittlinie der beiden Projektionsebenen oder vielmehr als $y$-Achse in der beibehaltenen Projektionsebene eine Parallele zu der Geraden, die die beiden Projektionspunkte verbindet.

Gefragt ist, wie das Stück der Kugeloberfläche, das in der ersten Projektion als Kreis erschien, in der zweiten erscheinen wird.

Koordinaten des Mittelpunkts dieses Flächenstücks (Radius $= 1$):

\begin{center}
\begin{tabular}{ccccc}
in & der & ersten & Projektion & $x,y,z$\\
,, & ,, & zweiten & ,, & $x',y',z'$.
\end{tabular}
\end{center}

Der Drehungswinkel $= \lambda$.

Es ist
\[
x' = x
\]
\[
\BareRoot{(1-xx)} = \rho
\]
\[
\BareRoot{(1-x'x')} = \rho
\]
\[
\frac{y}{\rho} = \sin u \qquad \frac{y'}{\rho} = \sin(u+\lambda)
\]
\[
dy' = \frac{z'}{z}\cdot dy - \frac{x\sin\lambda}{z}\cdot dx.
\]
